\apendice{Anexo de sostenibilización curricular}

\section{Introducción}
Hablar de sostenibilidad en el ámbito de la ingeniería informática no se limita únicamente a medir el consumo energético de los servidores o el hardware. Implica diseñar soluciones eficientes y útiles que aporten un valor real al entorno social y económico \cite{crue_sostenibilidad_2012}. En este TFG, enfocado en automatizar y analizar los datos turísticos de la provincia de Burgos, la sostenibilidad se ha integrado desde el propio diseño del software. Optimizar el código para no desperdiciar recursos de computación y visibilizar zonas con menor impacto turístico permite alinear este proyecto con los Objetivos de Desarrollo Sostenible (ODS) de la Agenda 2030 \cite{onu_agenda2030}. De este modo, se demuestra cómo las herramientas de datos pueden ser un motor de cambio en el desarrollo local.

\subsection{ODS 8: Trabajo decente y crecimiento económico - Impulso frente a la despoblación rural}
La despoblación rural es uno de los problemas más graves que afronta la provincia de Burgos, afectando directamente al desarrollo económico de la región. El turismo sostenible es una vía clave para reactivar la economía de estos pequeños municipios. Al automatizar la recogida y el análisis de los puntos de interés (POIs) y sus reseñas en localidades con pocos habitantes, el sistema ofrece a las administraciones una información que, de forma manual, sería inviable o demasiado costosa de obtener. Conocer el comportamiento y las opiniones de los visitantes en estas zonas permite a los gestores públicos diseñar campañas de promoción mucho más precisas y eficaces \cite{onu_ods8}. Esto ayuda a repartir los flujos turísticos fuera de las rutas habituales, impulsando los negocios locales, el comercio de proximidad y el empleo en municipios que sufren el declive demográfico.

\subsection{ODS 9: Industria, innovación e infraestructura - Eficiencia tecnológica y relación coste-beneficio}
Desde la perspectiva puramente técnica de la ingeniería de software, la sostenibilidad está ligada a la eficiencia de los algoritmos y al uso inteligente de la infraestructura cloud. El pipeline de datos se ha diseñado para funcionar estrictamente bajo demanda. Esto evita mantener servidores encendidos de forma continua y reduce el consumo energético innecesario durante los periodos en los que el scraper no necesita ejecutarse. La elección de una base de datos serverless como Google BigQuery optimiza el impacto ambiental al delegar la computación en una infraestructura global compartida y altamente eficiente. Además, el proyecto destaca por su alta viabilidad económica: con un coste de ejecución muy ajustado —gracias al aprovechamiento de los niveles gratuitos de GCP y al uso bajo demanda de Apify— se consigue procesar un volumen masivo de datos útiles para la gestión pública \cite{onu_ods9}. Innovar con arquitecturas eficientes demuestra que es posible desplegar soluciones estables a largo plazo sin sobrecargar los presupuestos institucionales.

\subsection{ODS 11: Ciudades y comunidades sostenibles - Preservación y visibilización del patrimonio natural}
El ODS 11 marca la importancia de proteger y poner en valor el patrimonio cultural y natural del mundo \cite{onu_ods11}. La provincia de Burgos posee una gran cantidad de espacios naturales, rutas y parajes que muchas veces carecen de la presencia digital necesaria para analizar su impacto real. Por ello, el pipeline incluye categorías específicas dedicadas al turismo rural y de naturaleza. Al recopilar y estructurar las opiniones de los usuarios sobre estos enclaves, el cuadro de mando en Looker Studio permite a los responsables de turismo evaluar cómo perciben los visitantes la conservación del entorno. Esta información ayuda a detectar problemas de masificación, planificar el mantenimiento de las infraestructuras o identificar qué áreas naturales despiertan más interés, promoviendo un modelo de turismo respetuoso con el medio ambiente y la biodiversidad local.

\section{Conclusión y reflexión personal}
El desarrollo de este TFG me ha servido para comprender que el desarrollo de software conlleva un impacto social directo que los ingenieros debemos asumir. Integrar estas competencias a lo largo del trabajo me ha enseñado que optimizar el código para reducir el tiempo de computación, usar estándares oficiales como el código de municipio del INE para evitar datos basura o plantear un almacenamiento en la nube eficiente, no son solo decisiones de ingeniería para ahorrar costes, son prácticas que minimizan la huella digital y aportan herramientas tecnológicas reales para combatir problemas tan complejos como la desigualdad territorial en nuestra provincia.